\chapter{Introduzione al MATLAB}

\begin{figure}[h]
    \centering
    
\includegraphics[scale = 0.1]{Matlab_Logo.png}
\end{figure}  

\newpage 

\section{Matlab ed i suoi utilizzi}

\footnote{ 
    Slide del prof | 2 - Introduzione al MATLAB | pag 2 - 3\\
    Appunti | 2025-09-30 | pag 3 - 5
}

\begin{tcolorbox}
    Il corso di DSP non è un corso di Matlab, ma useremo il programma MATLAB per imparare dei concetti relativi al DSP. \newline 

    Essendo studenti della mitica UNIVPM e se non avete mai scaricato MATLAB, 
    potete seguire la seguente guida \url{https://www.ict.univpm.it/node/236}. \newline 

    Visto che ci possono essere alcune differenze nell'uso dei comandi (cioè lo stesso comando tra una versione e l'altra può cambiare lievemente funzione), 
    in questo anno accademico 2025/2026, 
    il professore sta utilizzando la versione di MATLAB 2022b, invece nell'aula CESMI, cioè l'aula in cui si svolgerà l'esame, 
    è installata la versione 2023a. \newline 

    Il mio consiglio è quello di installare tutte e due le versioni con i due toolbox che serviranno all'esame, che sono: 
    
    \begin{itemize}
        \item Signal Processing Toolbox, link della relativa documentazione \\\url{https://it.mathworks.com/help/signal/index.html}
        \item DSP System Toolbox, link della relativa documentazione \\\url{https://it.mathworks.com/help/dsp/index.html}  
    \end{itemize}
    
    Vi consiglio di scaricare ADESSO le versioni di Matlab che vi servono perchè a Maggio 2025 c'è stato un attacco informatico che ha disattivato tutti i servizi di Mathworks per un mese intero\\ 
    \url{https://it.mathworks.com/company/trust-center/may-2025-ransomware-incident.html} \newline 
    
    Non voglio che, nel mentre che vi state preparando per l'esame, stiate con l'ansia che questo attacco possa capitare di nuovo e non potervi esercitare con questo programma. \newline 

    Se non avete mai programmato in vita vostra (non siamo nati tutti imparati, sennò non saremmo andati all'università), 
    la documentazione è quel file (ma generalmente è un sito) in cui potete approfondire meglio quello specifico argomento. \newline 

    Inoltre, Matlab mette a disposizione anche dei video come il seguente sia nel proprio sito web \\ 
    \url{https://it.mathworks.com/videos/introduction-to-signal-processing-apps-in-matlab-1580188288390.html}, 
    che nel suo canale su Youtube \url{https://www.youtube.com/@MATLAB/featured}. \newline 

    Se siete alle prime armi, o anche per un semplice ripasso in compagnia, vi consiglio il seguente video (in cui si utilizza Matlab 2020b ma le basi, come scritto precedentemente, più o meno sono le stesse): \newline 

    MATLAB Crash Course for Beginners by freeCodeCamp.org \\ 
    \url{https://www.youtube.com/watch?v=7f50sQYjNRA} \newline 

    \begin{comment}
    Negli appunti cercherò anche, se possibile, di spiegare il codice passo passo (anche perchè all'esame ci verrà richiesto di scrivere uno script Matlab). \newline 
    \end{comment}
    
    Ora basta con le doverose introduzioni, proseguiamo. 
\end{tcolorbox}

Matlab è un ambiente di sviluppo per algoritmi orientato all'elaborazione di matrici. \newline 

È un linguaggio ad "alte" prestazioni per problemi computazionali. \newline 

Matlab è disponibile su diverse piattaforme, MacOS, Linux, Windows (e anche sul cloud) e, ogni anno, 
vengono rilasciate due versioni: nella prima metà dell'anno esce la versione a e nella seconda metà dell'anno esce la versione b. \newline 

Generalmente è consigliato il Matlab per una prototipazione di un progetto e non da utilizzarlo per un progetto finale perchè è un linguaggio molto lento, 
cioè per svolgere un'istruzione ci mette tanto tempo e tante risorse (consiglio il Matlab per persone che devono sostituire il termosifone con il proprio pc in cameretta d'inverno): 
generalmente vengono utilizzati, successivamente, altri linguaggi più efficienti per finalizzare il progetto, come ad esempio il C/C++. \newline 

Oltre al Matlab si possono installare dei Toolbox, che sono dei programmi accessori al Matlab (un po' come delle espansioni nei giochi da tavola, ma in questo caso per Matlab). \newline 

Di seguito i Toolbox utili per il corso di DSP: 

    \begin{itemize}
        \item Signal Processing Toolbox \url{https://it.mathworks.com/products/signal.html}
        \item DSP System Toolbox \url{https://it.mathworks.com/products/dsp-system.html}  
    \end{itemize}

\newpage 

\section{Il sistema Matlab, l'IDE ed i Data Types}

\footnote{ 
    Slide del prof | 2 - Introduzione al MATLAB | pag 6 - 9\\
    Appunti | 2025-10-01 | pag 1 
}

Il sistema Matlab offre diverse possibilità di funzioni, ma noi ci concentreremo sulle principali. \newline 

Di seguito uno screenshot all'interno di Matlab, con delle note: 

\begin{figure}[h]
    \centering
    \includegraphics[scale = 0.3]{Schermata Matlab.png}
\end{figure} 

La command history è la finestra in cui potrete visualizzare i comandi più recenti. \newline  

Il workspace è la finestra in cui vengono memorizzati i dati più recenti. \newline 

La set path è quella dialog box dove indica la directory (cioè la cartella) in cui si trova il Matlab. \newline 

La command window è quella finestra dell'IDE per eseguire i comandi di Matlab. \newline 

Di seguito alcuni comandi più utilizzati, con gli usi più frequenti e con il link alla documentazione: 

\begin{itemize}
    \item save, per salvare un elemento da Matlab su un file esterno \\ \url{https://it.mathworks.com/help/matlab/ref/save.html} 
    \item load, per caricare un file esterno come elemento su Matlab \\ \url{https://it.mathworks.com/help/matlab/ref/load.html} 
    \item help, per visualizzare come può essere utilizzato un comando \\ \url{https://it.mathworks.com/help/matlab/ref/help.html} 
    \item clear, per cancellare delle variabili dal workspace \\ \url{https://it.mathworks.com/help/matlab/ref/clear.html} 
    \item whos, per avere delle informazioni dettagliate su un file \\ \url{https://it.mathworks.com/help/matlab/ref/whos.html} 
    \item pwd, per indicare la cartella in cui ci si trova \\ \url{https://it.mathworks.com/help/matlab/ref/pwd.html}  
    \item cd, per cambiare cartella \\ \url{https://it.mathworks.com/help/matlab/ref/cd.html} 
    \item dir, per elencare i file presenti nella cartella \\ \url{https://it.mathworks.com/help/matlab/ref/dir.html}
\end{itemize}

\begin{tcolorbox}
    Scrivo che generalmente quel tipo di comando viene utilizzato in questa maniera perchè, 
    se leggete la documentazione, lo stesso comando, in base alla versione ed ai parametri che inserite, 
    comporta ad esiti differenti. \newline 

    Durante l'esame scritto non potrete andare su internet per cercare di scrivere il proprio script: 
    ecco perchè il comando help e la documentazione interna saranno i vostri alleati. \newline 

    Se possibile, imparate già da adesso a leggerla ed a capirla, in modo che all'esame non perdiate del tempo: sarà solo un veloce ripasso leggere la documentazione.  
\end{tcolorbox} 

In Matlab i dati possono essere di diversi tipi. \newline 

Di seguito elenco i principali: 

\begin{itemize}
    \item Array multidimensionali classici, cioè scalari, vettori, matrici (come quelli che hai studiato nel corso di Algebra Lineare) 
    \item Array di celle, cioè array multidimensionali i cui elementi sono copie di altri array (un po' come le matrioske russe) 
    \item Caratteri, cioè degli array ottenuti scrivendo il testo tra apici 
    \item Strutture, cioè array multidimensionali accessibili attraverso designatori di campo  
\end{itemize}


\begin{tcolorbox}
    Per array si intende, dall'inglese, vettore. \newline 

    In questo corso, verranno utilizzati principalmente gli array multidimensionali e gli array di celle. \newline 

    Altro particolare di Matlab è che è un "linguaggio di programmazione" case-sensitive, cioè un linguaggio in cui le lettere maiuscole e minuscole distinguono due tipi di elementi. \newline 

    Ad esempio: pippo e Pippo sono due variabili differenti in Matlab. \newline 

    Inoltre, nel Matlab, gli indici partono da 1 e non da 0, come ad esempio nel C
\end{tcolorbox} 

\newpage 

\section{Operazioni su matrici e vettori}

\footnote{ 
    Slide del prof | 2 - Introduzione al MATLAB | pag 10 - 13\\
    Appunti | 2025-10-01 | pag 1 - 3 
}

Matlab è un linguaggio di programmazione orientato ai vettori e alle matrici, quindi vediamo un po' di codici di esempi per generarli. \newline 

\begin{tcolorbox}
Se volete approfondire le operazioni che si possono fare con le matrici su Matlab\\  
\url{https://it.mathworks.com/help/matlab/math/matrices-in-the-matlab-environment.html}    
\end{tcolorbox}

\textbf{Generazione di matrici} \newline 

Di seguito una generazione di una matrice M nella command window: \newline 

\begin{figure}[h]
    \centering
    \includegraphics[scale = 0.4]{Matrice M codice.png}
\end{figure} 

\begin{tcolorbox}
In base a come abbiamo scritto il comando, 
la variabile M è una matrice 3x3, cioè 3 righe e 3 colonne. \newline 

Per generare il vettore M bisogna inserire i valori all'interno delle parentesi quadre, cioè $[]$. \newline 

Il punto e virgola (;) permette di scrivere la riga successiva.     
\end{tcolorbox} 

\textbf{Estrazione di una sottomatrice} \newline 

Di seguito l'estrazione di una sotto-matrice da command window in cui è disponibile la matrice M nel workspace: 

\begin{figure}[h]
    \centering
    \includegraphics[scale = 0.4]{Sotto-matrice Matlab comando.png}
\end{figure} 

\begin{tcolorbox}
    Con quel comando indicato nella command window, 
    si intende che si vuole prendere la sottomatrice di M che va dalla prima riga alla seconda riga e dalla seconda colonna alla terza colonna. \newline 
    
    In questo caso, si utilizzano i due punti (:) per dire da quale riga/colonna si vuole partire a sinistra dei due punti, a quella riga/colonna indicato a destra dei due punti.  
\end{tcolorbox} 

\newpage 

\textbf{Operatore colon ":"} \newline 

Di seguito un esempio dell'operatore colon ":" usato frequentemente nei comandi: 

\begin{figure}[h]
    \centering
    \includegraphics[scale = 0.4]{Vettore con operatore colon.png}
\end{figure} 

\begin{tcolorbox}
    L'operatore colon ":" è molto utile perchè piuttosto che scrivere: 
    
    \begin{lstlisting}[
    frame=single,
    style=Matlab-Pyglike]
a = [0 0.1 0.2 0.3 0.4 0.5 0.6 0.7 0.8 0.9 1]

    \end{lstlisting}

    Il comando: 
    \begin{lstlisting}[
    frame=single,
    style=Matlab-Pyglike]
a = [0: 0.1 : 1]

    \end{lstlisting}
    è molto più rapido e agevole. \newline 

    In altre parole il comando esprime di creare un vettore che parte dal valore iniziale 0 fino al valore finale 1 a passo di 0.1. \newline 

    Il passo può essere anche negativo. 
\end{tcolorbox} 

\textbf{Funzioni matriciali} \newline 

\begin{tcolorbox}
Se volete approfondire le operazioni che si possono fare con le matrici su Matlab \\ 
\url{https://it.mathworks.com/help/matlab/math/matrices-in-the-matlab-environment.html}    
\end{tcolorbox} 

Esistono molte funzioni matriciali già implementate in Matlab. \newline 

Ad esempio, il comando: 

\begin{lstlisting}[
    frame=single,
    style=Matlab-Pyglike]
M_trasposta = M'  
\end{lstlisting}

implementa la funzione matriciale: 

{
    \Large 
    \begin{equation}
        M_{trasposta} = M^{T}
    \end{equation}
}

dove M è una matrice e con T si intende la sua trasposta. \newline 

\begin{tcolorbox}
    Se non ti ricordi cosa significa fare la matrice trasposta: \newline 

    Matrice Trasposta by Wikipedia \\ 
    \url{https://it.wikipedia.org/wiki/Matrice_trasposta}
\end{tcolorbox} 

Il comando sum in esempio: 

\begin{lstlisting}[
frame=single,
numbers=left,
style=Matlab-Pyglike]
vect = [1 0 -1 2 1 4 2]; 
somma_elementi_vettore = sum(vect)
\end{lstlisting}

somma tutti gli elementi contenuti nel vettore vect, cioè implementa la seguente funzione matematica: 

{
    \Large 
    \begin{equation}
        \begin{split}
    vect = 1 + 0 -1 +2 +1 +4 + 2        
        \end{split}
    \end{equation}
}

Il comando diag in esempio: 

\begin{lstlisting}[
    frame=single,
    style=Matlab-Pyglike]
diag(M)  
\end{lstlisting}

calcola elenca gli elementi di una diagonale di una matrice. \newline 

\begin{tcolorbox}
    Per apprfondire il comando diag: \\\url{https://it.mathworks.com/help/matlab/ref/diag.html} \newline 

    Se non ti ricordi cosa significa la diagonale di una matrice: \newline 

    Matrice diagonale by Wikipedia \\  
    \url{https://it.wikipedia.org/wiki/Matrice_diagonale}
\end{tcolorbox}

Il comando eig in esempio: 

\begin{lstlisting}[
    frame=single,
    style=Matlab-Pyglike]
eig(M)  
\end{lstlisting}

calcola gli autovalori della matrice M. \newline 

\begin{tcolorbox}

    Per approfondire il comando eig: \\\url{https://it.mathworks.com/help/matlab/ref/eig.html}\newline 

    Se non ti ricordi cosa significa la diagonale di una matrice: \newline 

   Eigenvectors and eigenvalues | Chapter 14, Essence of linear algebra by 3Blue1Brown \\ 
   \url{https://youtu.be/PFDu9oVAE-g?si=wJ0_hB40hehyk4Qq}
\end{tcolorbox} 

Il comando per il vettore M2: 

\begin{lstlisting}[
frame=single,
numbers=left,
style=Matlab-Pyglike]
b = [5 5 5]; 
M2 = [M; b]
\end{lstlisting}

permette di concatenare due o più matrici (in questo caso la matrice M2 è concatenata alla matrice M e alla matrice b). \newline 

\begin{tcolorbox}
    Per visualizzare un esempio di concatenazione di matrici: \newline 

    Matrix Concatenation in MATLAB by thesalmonification \\ 
    \url{https://youtu.be/TAIetkMXmLs?si=Qf_-3bJuiA75J65_}
\end{tcolorbox}

Oltre alle funzioni matriciali, è possibile definire delle matrici multidimensionali, 
i quali vengono trattate come strutture bidimensionali. \newline 

Ad esempio i seguenti comandi: 

\begin{lstlisting}[
frame=single,
numbers=left,
style=Matlab-Pyglike]
A = randn(2, 3, 2); 
A(:, :, 1)
\end{lstlisting}

permettono di creare una matrice multidimensionale con dimensione 2x3x2 con valori casuali. \newline 

\begin{tcolorbox}
   Per approfondire gli array multidimensionali: \\\url{https://it.mathworks.com/help/matlab/math/multidimensional-arrays.html}
\end{tcolorbox}

\newpage 

\newpage 

Oltre a questo, è possibile svolgere delle operazioni matriciali come le seguenti: 

\begin{figure}[h]
    \centering
    \includegraphics[scale = 0.4]{Operazioni matriciali Matlab.png}
\end{figure} 

Inoltre, utilizzando il punto prima dell'operatore, 
si può effettuare un'operazione elemento per elemento (cioè element-wise) tra i vettori e matrici coinvolti. \newline 

\begin{tcolorbox}
    Un approfondimento su questo argomento: \\\url{https://it.mathworks.com/help/matlab/matlab_prog/array-vs-matrix-operations.html}
\end{tcolorbox}

Di seguito un esempio di comandi utilizzati comunemente su Matlab: 

\begin{figure}[h]
    \centering
    \includegraphics[scale = 0.55]{Comandi matlab utili.png}
\end{figure} 


Come scritto nei capitoli precedenti, il tempo nel DSP non esiste perchè siamo nell'ambito dei vettori. \newline 

Generalmente assoceremo il tempo in base a dove si trova nel vettore il valore (cioè l'indice del vettore sarà per noi l'istante temporale). \newline 

I numeri complessi esistono in Matlab: la parte immaginaria viene indicata o con la variabile i o con la variabile j. \newline 

Ad esempio se abbiamo il numero: 

{
    \Large 
    \begin{equation}
        \overrightarrow{A} = 3 + 3 \jmath
    \end{equation}
}

in Matlab lo possiamo scrivere come: 

\begin{lstlisting}[
    frame=single,
    style=Matlab-Pyglike]
A = 3 + 3i
\end{lstlisting}

oppure: 

\begin{lstlisting}[
    frame=single,
    style=Matlab-Pyglike]
A = 3 + 3j
\end{lstlisting}

\begin{comment}
\begin{tcolorbox}[colback=green!5!white,colframe=green!75!black]
TO-DO: Spiegazione codice Esempio 2.1 
\end{tcolorbox} 
\end{comment}

\newpage 

\section{Plotting}

\footnote{ 
    Slide del prof | 2 - Introduzione al MATLAB | pag 14 - 15\\
    Appunti | 2025-10-01 | pag 3
}

In Matlab, ci sono diversi modi di graficare i vettori, che sono: 

\begin{itemize}
    \item plot, in cui vengono interpolati (cioè collegati) i punti in "modo" continuo \\\url{https://it.mathworks.com/help/matlab/ref/plot.html} 
    \item stem, in cui viene data una rappresentazione ad impulsi \\\url{https://it.mathworks.com/help/matlab/ref/stem.html}
    \item bar, in cui viene data una rappresentazione con i diagrammi a barre \\\url{https://it.mathworks.com/help/matlab/ref/bar.html} 
\end{itemize}

\begin{tcolorbox}
    Se vuoi sapere tutti i modi per graficare un vettore: \\\url{https://it.mathworks.com/help/matlab/creating_plots/types-of-matlab-plots.html}
\end{tcolorbox}

Se si invoca un solo vettore, esso verrà graficato rispetto all'indice. \newline 

Ad esempio dato il vettore con valori casuali A: 

\begin{lstlisting}[
frame=single,
numbers=left,
style=Matlab-Pyglike]
A = randn(1, 10); 
plot(A)
\end{lstlisting}

avremo il seguente caso: 

\begin{figure}[h]
    \centering
    \includegraphics[scale = 0.7]{plot di esempio con un vettore.png}
\end{figure} 

dove, nell'asse delle x o chiamato anche l'asse delle ascisse, viene utilizzato l'indice del vettore. \newline 

\newpage 

Se invece si invocano due vettori, il primo è utilizzato nell'asse delle x, cioè nell'asse delle ascisse, ed il secondo per l'asse y, cioè quello delle ordinate. \newline 

Ad esempio, dati i seguenti vettori A e B in cui B ha valori casuali: 

\begin{lstlisting}[
frame=single,
numbers=left,
style=Matlab-Pyglike]
A = [0: 0.1 : 1]
B = randn(1, 10); 
plot(A, B)
\end{lstlisting}


\begin{figure}[h]
    \centering
    \includegraphics[scale = 0.7]{plot di esempio con due vettori.png}
\end{figure} 

Quindi, in questo caso ad esempio, il vettore A può essere utilizzato per indicare "il tempo", 
invece il vettore B per indicare i valori istantanei della velocità di un maratoneta. \newline 

Di default, Matlab dirige l'output dei comandi grafici sulla finestra figure più recente (ad esempio su Windows si apre un'altra finestra oltre a quella dell'IDE di Matlab chiamata, ad esempio, Figure 1). \newline 

I seguenti comandi sono utili per gestire più figure in Matlab: 

\begin{itemize}
    \item figure, crea una nuova finestra "figure" e vi dirige l'output 
    \item figure(N), dirige l'output sulla finestra N 
    \item hold on, abilita la sovrapposizione dei grafici sull'ultima figure attivata 
    \item hold off, disabilita la sovrapposizione 
    \item sublot(r, c, i), suddivide la finestra in una matrice di sottografici di r righe, c colonne e direzione l'output sul sottografico i-esimo (contando per righe)
\end{itemize}

\begin{tcolorbox}
    Per approfondire i comandi figure, hold, subplot: 

    \begin{itemize}
        \item figure \url{https://it.mathworks.com/help/matlab/ref/figure.html} 
        \item hold \url{https://it.mathworks.com/help/matlab/ref/hold.html} 
        \item subplot \url{https://it.mathworks.com/help/matlab/ref/subplot.html}
    \end{itemize} 

\end{tcolorbox}

\begin{comment}
\begin{tcolorbox}[colback=green!5!white,colframe=green!75!black]
TO-DO: Spiegazione codice Esempio 2.2 
\end{tcolorbox} 
\end{comment}

\newpage 

\begin{comment}
\section{Esercizi di esempio}

\footnote{ 
    Slide del prof | 2 - Introduzione al MATLAB | pag 16 - 17\\
    Appunti | 2025-10-01 | pag 3
}

\begin{tcolorbox}[colback=green!5!white,colframe=green!75!black]
TO-DO: Spiegazione codice Esempio 2.1, 2.2, 2.3, 2.4 
\end{tcolorbox} 

\newpage 

\end{comment}

\section{Matlab Programming}

\footnote{ 
    Slide del prof | 2 - Introduzione al MATLAB | pag 18\\
    Appunti mancanti | 2025-10-07 
}

È possibile realizzare programmi ad-hoc tramite il linguaggio di programmazione Matlab. \newline 

Il linguaggio Matlab, oltre alle semplici assegnazioni di variabili e ad operatori matematici, comprende anche "statements" più evoluti come: 

\begin{itemize}
    \item if, elseif, else, eseguono  delle operazione se le condizioni sono vere \\\url{https://it.mathworks.com/help/matlab/ref/if.html} 
    \item switch, case, otherwise, eseguono uno tra i gruppi di operazioni \\\url{https://it.mathworks.com/help/matlab/ref/switch.html} 
    \item for, foor-loop che si ripete un numero di volte prefissato \\\url{https://it.mathworks.com/help/matlab/ref/for.html} 
    \item while, while-loop che si ripete se la condizione rimane vera \\\url{https://it.mathworks.com/help/matlab/ref/while.html} 
    \item break, termina un'esecuzione di un while-loop o di un for-loop \\\url{https://it.mathworks.com/help/matlab/ref/break.html}
\end{itemize}

Esistono due possibilità per programmare in Matlab: 

\begin{itemize}
    \item script 
    \item funzioni 
\end{itemize} 

\newpage 

\subsection{Script file}

\footnote{ 
    Slide del prof | 2 - Introduzione al MATLAB | pag 19 - 20\\
    Appunti mancanti | 2025-10-07 
}

Il file script è un file che contiene una sequenza di comandi Matlab, generalmente con estensione .m. \newline 

È bene non assegnare a nessun script il nome delle funzioni predefinite o il nome di una variabile usata. \newline 

\begin{comment}
\begin{tcolorbox}[colback=green!5!white,colframe=green!75!black]
TO-DO: Spiegazione codice Esempio 2.5
\end{tcolorbox} 
\end{comment}

\newpage 

\subsection{Funzioni}

\footnote{ 
    Slide del prof | 2 - Introduzione al MATLAB | pag 21\\
    Appunti mancanti | 2025-10-07 
}

Una funzione si crea come uno script: ha estensione .m, deve avere lo stesso nome della funzione, 
ma, a differenza dello script, può restituire dei risultati i cui parametri avvengono per valore. \newline 

Il passaggio per valore implica che le variabili definite all'interno della funziona hanno visibilità locale e non possono essere riferite dalla shell. \newline 

\begin{tcolorbox}
    I concetti di visibilità locale e globale sono dei concetti molto importanti nella programmazione. \newline 

    Per approfondire: \\\url{https://it.wikipedia.org/wiki/Visibilit%C3%A0_(programmazione)} \newline 

    Per quanto riguarda il passaggio del valore in una funzione, lo si ha anche in altri linguaggi di programmazione, come ad esempio il C e il C++. \newline 

    Per approfondire la differenza tra passaggio per valore e passaggio per riferimento: 
    
    \begin{itemize}
        \item Passaggio per valore \\ \url{https://www.ibm.com/docs/it/i/7.5.0?topic=calls-pass-by-value}
        \item Passaggio per riferimento \\ \url{https://www.ibm.com/docs/it/i/7.5.0?topic=calls-pass-by-reference}
    \end{itemize}
    
\end{tcolorbox}

In Matlab esiste un tipo di variabile interessante chiamata "Persistent variable", 
che è una variabile locale che vista come globale solamente dalla funzione dove è stata definita. \newline 

\begin{tcolorbox}
    Nella documentazione di persistent sono disponibili degli esempi più dettagliati per approfondire il tema: \\ 
    \url{https://it.mathworks.com/help/matlab/ref/persistent.html}
\end{tcolorbox}

\newpage 

\section{Convoluzione di due sequenze}

\footnote{ 
    Slide del prof | 2 - Introduzione al MATLAB | pag 22\\
    Appunti mancanti | 2025-10-07 \\ 
    Latex | 2025-10-13 Sera | pag 59 - 60
}

Come studiato dal capitolo precedente, sappiamo bene che 
la convoluzione tra la sequenza discreta di ingresso ad un sistema discreto LTI $x \left[ n\right]$ e la risposta impulsiva $h \left[ n\right]$, 
è definita come: 

{
    \Large 
    \begin{equation}
        y \left[ n\right] 
        = 
        \sum_{k = - \infty}^{+ \infty}
        x \left[ k\right]
        \cdot 
        h \left[ n - k\right]
    \end{equation}
}

In Matlab esiste una funzione pre-implementata per il calcolo della convoluzione tra due sequenze ed è la seguente: 

{
    \Large 
    \begin{lstlisting}[
    frame=single,
    style=Matlab-Pyglike]
y = conv(x, h)
\end{lstlisting}
}

Ma, come si legge bene dalla documentazione di conv \url{https://it.mathworks.com/help/matlab/ref/conv.html}, 
essa produce un vettore y di lunghezza pari a: 

{
    \Large 
    \begin{lstlisting}[
    frame=single,
    style=Matlab-Pyglike]
length(x) + length(h) - 1
\end{lstlisting}
} 

mentre noi vorremmo che la lunghezza del vettore y sia: 

{
    \Large 
    \begin{lstlisting}[
    frame=single,
    style=Matlab-Pyglike]
length(x) + length(h) 
\end{lstlisting}
} 

Per questo motivo che, in questo corso, useremo la funzione my\_conv.m .\newline 

\begin{comment}
\begin{tcolorbox}[colback=green!5!white,colframe=green!75!black]
TO-DO: Spiegazione codice Esempio 2.3 my\_conv.m
\end{tcolorbox} 
\end{comment}

Come studiato nel capitolo precedente \ref{subsec: Proprietà di campionamento} \nameref{subsec: Proprietà di campionamento} 
una qualsiasi sequenza $x \left[ n \right]$ può essere espressa come somma di impulsi unitari 
opportunamente ritardati e amplificati. \newline 

Da un punto di vista matematico, possiamo scrivere la sequenza $x \left[ n \right]$ come: 

{
    \Large 
    \begin{equation}
        x \left[ n \right] 
        = 
        \sum_{k = - \infty}^{+ \infty }
        x \left[ k \right] \cdot \delta \left[ n - k \right]
    \end{equation}
}

dove il singolo impulso unitari $\delta \left[ n\right]$ prende il nome di Delta di Kronocker. \newline 

La proprietà di campionamento delle sequenze discrete appena descritta, permette di assegnare uno specifico significato operativo all'operazione di convoluzione. \newline 

\begin{comment}
\begin{tcolorbox}[colback=green!5!white,colframe=green!75!black]
TO-DO: Spiegazione codice Esempio 2.4 
\end{tcolorbox}
\end{comment}


Tale significato operativo può essere analogamente essere messo in luce considerando la convoluzione come un'operazione matriciale. \newline 

L'operazione di convoluzione tra due sequenze può essere vista come il prodotto matriciale tra una delle due sequenze e la matrice di convoluzione, opportunamente costruita a partire dalla sequenza rimanente. \newline 

La matrice di convoluzione $\left[ H\right]$ è una matrice di Toepliz rettangolare di dimensione N * p dove: 

\begin{itemize}
    \item N è la lunghezza del prodotto di convoluzione 
    \item p è la lunghezza della sequenza x che non genera la matrice 
\end{itemize}

\begin{tcolorbox}

Possiamo dire che una matrice di Toepliz utile per la proprietà di campionamento nei sistemi discreti 
è ottenuta utilizzando la Delta di Kronocher, 
che appunto alla fine compone una matrice identità. \newline 

Ma, non è detto che la matrice di Toepliz debba per forza essere simmetrica (approfondiremo meglio alla magistrale). \newline 

Per approfondire cosa è una matrice di Toepliz e la Delta di Kronocher: 

\begin{itemize}
    \item \url{https://it.wikipedia.org/wiki/Matrice_di_Toeplitz} 
    \item \url{https://it.wikipedia.org/wiki/Delta_di_Kronecker}
\end{itemize}

\end{tcolorbox}

La matrice è ottenuta shiftatndo nel tempo la sequenza base h su ogni colonna. \newline 

\begin{comment}
Per esempio: 

\begin{tcolorbox}[colback=green!5!white,colframe=green!75!black]
TO-DO: Fare un esempio di una convoluzione tra due sequenze utilizzando la matrice di Toeplitz  
\end{tcolorbox} 
\end{comment}

Questa interpretazione è molto utile ai fini computazionali. \newline 

Il comando Matlab da utilizzare per generare la matrice di convoluzione è convmtx. \newline 


\begin{tcolorbox}
    Dalla documentazione di convmtx \\\url{https://it.mathworks.com/help/signal/ref/convmtx.html} \newline 

    "Computing a convolution using conv when the signals are vectors is generally more efficient than using convmtx. For multichannel signals, convmtx might be more efficient." \newline 

    convmtx ha lo stesso problema del comando conv in Matlab, 
    cioè della sua grandezza che non rispecchia quella ideale come noi la intendiamo nel corso: 
    ecco perchè dovremo costrutire la nostra my\_convmtx.m 

\end{tcolorbox} 

\begin{comment}
\begin{tcolorbox}[colback=green!5!white,colframe=green!75!black]
TO-DO: TO-DO: Spiegazione codice Esempio 2.5, 2.6, 2.7 
\end{tcolorbox} 
\end{comment}

\newpage 
. 

\newpage 